\documentclass[11pt]{article}
\usepackage{hw-solutions-common}

\begin{document}

\hwsolhead{Written HW 3}{Problem statements from Knight, \emph{Physics for Scientists and
Engineers} (5th ed., Global Edition), Chapter 23. $K = 8.99\times10^{9}$
N$\cdot$m$^2$/C$^2$, $\varepsilon_0 = 8.85\times10^{-12}$ C$^2$/(N$\cdot$m$^2$), and
$g = 9.8$ m/s$^2$ throughout.}

%============================================================
\problem{23.3}{2}{Two point charges, $q_A = +3.0$ nC and $q_B = -3.0$ nC, sit one above
the other with a 5.0 cm gap on each side of the point marked by the dot in Figure EX23.3:
$q_A$ is 5.0 cm above the dot's height, $q_B$ is 5.0 cm below it, and the dot itself is
5.0 cm to the side of the line joining the two charges. What are the strength and
direction of the electric field at the dot? Specify the direction as an angle cw from
horizontal.}

\model{Treat both charges as point charges. The field at the dot is the vector sum of
the field from $q_A$ alone and the field from $q_B$ alone.}

\visualize{Put the midpoint between the charges at the origin, $q_A$ at $(0, 5.0\text{
cm})$, $q_B$ at $(0, -5.0\text{ cm})$, and the dot at $(5.0\text{ cm}, 0)$. Both charges
are the same distance from the dot, along lines at $45^\circ$ to the vertical. $q_A$ is
positive, so its field points away from $q_A$, down and to the right. $q_B$ is negative,
so its field points toward $q_B$, down and to the left. The horizontal components point
opposite ways and will cancel; the vertical components both point down and will add.}

\solve{Each charge is a distance $r = \sqrt{(5.0\text{ cm})^2+(5.0\text{ cm})^2} =
7.07$ cm $= 0.0707$ m from the dot, so
\[
E_A = E_B = \frac{K|q|}{r^2} = \frac{(8.99\times10^{9})(3.0\times10^{-9})}{(0.0707\text{
m})^2} = 5.39\times10^{3}\text{ N/C}.
\]
Each field makes a $45^\circ$ angle with the vertical, so each has horizontal component
$E\cos45^\circ$ and vertical component $E\sin45^\circ$, both equal in magnitude since
$\sin45^\circ=\cos45^\circ$. The horizontal components are equal and opposite (one
points right, one points left), so they cancel exactly. The vertical components both
point downward and add:
\[
E_{\text{net}} = 2E_A\sin45^\circ = 2(5.39\times10^{3}\text{ N/C})(0.707) =
\boxed{7.6\times10^{3}\text{ N/C, straight down (}90^\circ\text{ cw from horizontal)}}.
\]}

\review{The horizontal cancellation isn't a coincidence: the dot is equidistant from
both charges and the charges have equal and opposite magnitude, so whatever field $q_A$
contributes sideways, $q_B$ contributes the exact opposite. Only the vertical
components, which both point the same way (from $+$ toward $-$), survive. That the net
field ends up parallel to the line joining the charges is a general feature of any
symmetric point on a dipole's perpendicular bisector.}

%============================================================
\problem{23.4}{2}{$q_A = +3.0$ nC sits at the same height as, and 5.0 cm to the left
of, the dot marked in Figure EX23.4. $q_B = -6.0$ nC sits 10 cm directly below $q_A$.
What are the strength and direction of the electric field at the dot? Specify the
direction as an angle cw from horizontal.}

\model{Treat both charges as point charges and add their fields as vectors.}

\visualize{Put $q_A$ at the origin, the dot at $(5.0\text{ cm}, 0)$, and $q_B$ at
$(0,-10\text{ cm})$. $q_A$ is positive and directly to the left of the dot, so its field
at the dot points straight in the $+x$ direction. $q_B$ is negative and down and to the
left of the dot, so its field at the dot points down and to the left, toward $q_B$.}

\solve{$q_A$ is a distance $r_A = 5.0$ cm $=0.050$ m from the dot:
\[
E_A = \frac{Kq_A}{r_A^2} = \frac{(8.99\times10^{9})(3.0\times10^{-9})}{(0.050)^2} =
1.08\times10^{4}\text{ N/C, in } +\hat{x}.
\]
$q_B$ is a distance $r_B = \sqrt{(5.0\text{ cm})^2+(10\text{ cm})^2} = 11.18$ cm
$=0.1118$ m from the dot, and the unit vector from the dot toward $q_B$ is
$(-5.0,-10)/11.18 = (-0.447,-0.894)$:
\[
E_B = \frac{K|q_B|}{r_B^2} = \frac{(8.99\times10^{9})(6.0\times10^{-9})}{(0.1118)^2} =
4.32\times10^{3}\text{ N/C},
\qquad
\vec{E}_B = (-1.93,-3.86)\times10^{3}\text{ N/C}.
\]
Adding components, $\vec{E}_A = (1.08\times10^4, 0)$ N/C, so
\[
\vec{E}_{\text{net}} = (8.86,-3.86)\times10^{3}\text{ N/C}
\;\Rightarrow\;
E_{\text{net}} = \sqrt{8.86^2+3.86^2}\times10^{3}\text{ N/C} = \boxed{9.7\times10^{3}\text{
N/C}},
\]
directed $\tan^{-1}(3.86/8.86) = \boxed{23^\circ\text{ below horizontal (cw)}}$, i.e.,
down and to the right.}

\review{$q_A$ is closer to the dot than $q_B$ (5.0 cm vs.\ 11.2 cm) and $q_A$'s field
points straight along $+\hat x$ with nothing to cancel it, so it's no surprise that
$q_A$ dominates the result: the net field is close to $E_A$'s value and tilted only
mildly toward $q_B$'s direction.}

%============================================================
\problem{23.8}{3}{A 10-cm-long thin glass rod uniformly charged to $+10$ nC and a
10-cm-long thin plastic rod uniformly charged to $-10$ nC are placed side by side, 4.0
cm apart. What are the electric field strengths $E_1$ to $E_3$ at distances 1.0 cm, 2.0
cm, and 3.0 cm from the glass rod, along the line connecting the midpoints of the two
rods?}

\model{Each rod is a finite uniformly charged line. On the line through a rod's own
midpoint, perpendicular to the rod, the field of a rod with charge $Q$ and length $L$ is
\[
E_{\text{rod}}(y) = \frac{KQ}{y\sqrt{y^2+(L/2)^2}},
\]
directed away from the rod if $Q>0$, toward it if $Q<0$. All three points lie on this
line for \emph{both} rods at once (they're between the rods), so the total field is the
vector sum of the two rods' contributions.}

\visualize{Put the glass rod's center at $x=0$ and the plastic rod's center at
$x=4.0$ cm. A point at distance $y_1$ from the glass rod (measured from $x=0$) is then a
distance $y_2 = 4.0\text{ cm} - y_1$ from the plastic rod. Glass is positive, so its
field pushes away from itself, toward the plastic rod ($+x$). Plastic is negative, so its
field pulls toward itself, which from a point between the rods is also toward $+x$. Both
fields point the same way at every point between the rods, so the strengths simply add.}

\solve{With $L=0.10$ m so $L/2=0.050$ m, and $Q=1.0\times10^{-8}$ C for each rod (same
magnitude, opposite sign):
\[
E_1:\ y_1=1.0\text{ cm},\ y_2=3.0\text{ cm}
\;\Rightarrow\; E_{\text{glass}}=1.76\times10^{5},\ E_{\text{plastic}}=5.14\times10^{4}
\;\Rightarrow\; E_1 = \boxed{2.3\times10^{5}\text{ N/C}}
\]
\[
E_2:\ y_1=y_2=2.0\text{ cm}
\;\Rightarrow\; E_{\text{glass}}=E_{\text{plastic}}=8.35\times10^{4}
\;\Rightarrow\; E_2 = \boxed{1.7\times10^{5}\text{ N/C}}
\]
\[
E_3:\ y_1=3.0\text{ cm},\ y_2=1.0\text{ cm}
\;\Rightarrow\; E_{\text{glass}}=5.14\times10^{4},\ E_{\text{plastic}}=1.76\times10^{5}
\;\Rightarrow\; E_3 = \boxed{2.3\times10^{5}\text{ N/C}}
\]
all three directed from the glass rod toward the plastic rod.}

\review{$E_1$ and $E_3$ come out identical. That has to be true: the point 1.0 cm from
glass (3.0 cm from plastic) is the mirror image of the point 3.0 cm from glass (1.0 cm
from plastic), with the roles of ``near rod'' and ``far rod'' simply swapped between two
rods that carry equal and opposite charge of equal magnitude. The field is \emph{smaller}
exactly in the middle than near either rod: moving toward a rod, that rod's own
contribution grows faster than the other rod's contribution shrinks.}

%============================================================
\problem{23.10}{2}{A small glass bead, charged to $+8.0$ nC, is in the plane that
bisects a thin, uniformly charged, 10-cm-long glass rod and is 4.0 cm from the rod's
center. The bead is repelled from the rod with a force of 720 $\mu$N. What is the total
charge on the rod?}

\model{The bead sits on the rod's perpendicular bisector, so its field there is
$E=KQ_{\text{rod}}/\big(y\sqrt{y^2+(L/2)^2}\big)$. The force on the bead is $F=qE$.}

\visualize{The bead is \emph{repelled}, so the rod's charge must have the same sign as
the bead's: the rod is positively charged.}

\solve{First find the field the bead sits in from the force it feels:
\[
E = \frac{F}{q} = \frac{7.20\times10^{-4}\text{ N}}{8.0\times10^{-9}\text{ C}} =
9.0\times10^{4}\text{ N/C}.
\]
With $y=0.040$ m and $L/2=0.050$ m, $y\sqrt{y^2+(L/2)^2} = (0.040)\sqrt{0.0016+0.0025} =
2.56\times10^{-3}$ m$^2$, so
\[
Q_{\text{rod}} = \frac{Ey\sqrt{y^2+(L/2)^2}}{K} =
\frac{(9.0\times10^{4})(2.56\times10^{-3})}{8.99\times10^{9}} =
\boxed{+26\text{ nC}}.
\]}

\review{26 nC on a 10-cm rod is a linear density of about 2.6 nC/cm, roughly comparable
to the point charges used elsewhere in this chapter's problems, so the answer is a
believable, everyday-lab-bench amount of charge rather than something wildly large or
small.}

%============================================================
\problem{23.12}{2}{The electric field 2.0 cm from a very long charged wire is 1000 N/C,
directed toward the wire. What is the charge (in nC) on a 5.0-cm-long segment of the
wire?}

\model{Far from its ends, a long charged wire produces the field of an infinite line of
charge, $E = \lambda/(2\pi\varepsilon_0 r) = 2K\lambda/r$, independent of the wire's
total length.}

\visualize{The field points \emph{toward} the wire, which is what a negative line charge
does (it attracts a positive test charge in toward itself). So the wire is negatively
charged.}

\solve{Solve for the linear charge density:
\[
\lambda = \frac{Er}{2K} = \frac{(1000\text{ N/C})(0.020\text{ m})}{2(8.99\times10^{9})} =
1.11\times10^{-9}\text{ C/m}.
\]
The charge on a 5.0-cm $=0.050$ m segment is then
\[
Q = \lambda L = (1.11\times10^{-9}\text{ C/m})(0.050\text{ m}) = 5.6\times10^{-11}\text{
C} = \boxed{-0.056\text{ nC}}
\]
(negative, since the wire itself is negatively charged).}

\review{The field formula for a line of charge doesn't care how long the wire is, only
how densely charged it is and how far away you're standing; the 5.0-cm length only
enters at the very last step, to convert the density into an actual amount of charge on
that particular stretch of wire.}

%============================================================
\problem{23.19}{3}{Two very large sheets of plastic face each other, separated by
distance $d$, as shown in Figure EX23.19. One sheet (on the left) has been rubbed to a
uniform surface charge density $\eta_1 = -\eta_0$; the other (on the right) has surface
charge density $\eta_2 = +3\eta_0$. What are the electric field vectors at points 1
(outside, to the left of both sheets), 2 (midway between the sheets), and 3 (outside, to
the right of both sheets)?}

\model{Each sheet, treated as an infinite plane of charge, produces a field of constant
magnitude $\eta/(2\varepsilon_0)$ on each side, independent of distance from that sheet:
pointing away from the sheet if $\eta>0$, toward it if $\eta<0$. The total field at any
point is the vector sum of both sheets' contributions.}

\visualize{Put the negative sheet at $x=0$ and the positive sheet at $x=d$, with $+x$
pointing to the right. Sheet 1 ($-\eta_0$) points its field toward itself everywhere: in
$+\hat x$ to its left (point 1) and in $-\hat x$ everywhere to its right (points 2 and
3). Sheet 2 ($+3\eta_0$) points its field away from itself everywhere: in $-\hat x$
everywhere to its left (points 1 and 2) and in $+\hat x$ to its right (point 3).}

\solve{Each sheet contributes magnitude $\eta_0/(2\varepsilon_0)$ (sheet 1) or
$3\eta_0/(2\varepsilon_0)$ (sheet 2). Adding the two contributions, with sign for
direction ($+\hat x$ positive):
\[
\text{Point 1 (left of both):}\quad
\frac{\eta_0}{2\varepsilon_0} - \frac{3\eta_0}{2\varepsilon_0} = -\frac{\eta_0}{\varepsilon_0}
\;\Rightarrow\;
\vec{E}_1 = \boxed{\dfrac{\eta_0}{\varepsilon_0}\text{, pointing left, away from the sheets}}
\]
\[
\text{Point 2 (between):}\quad
-\frac{\eta_0}{2\varepsilon_0} - \frac{3\eta_0}{2\varepsilon_0} = -\frac{2\eta_0}{\varepsilon_0}
\;\Rightarrow\;
\vec{E}_2 = \boxed{\dfrac{2\eta_0}{\varepsilon_0}\text{, pointing left, from sheet 2 toward sheet 1}}
\]
\[
\text{Point 3 (right of both):}\quad
-\frac{\eta_0}{2\varepsilon_0} + \frac{3\eta_0}{2\varepsilon_0} = +\frac{\eta_0}{\varepsilon_0}
\;\Rightarrow\;
\vec{E}_3 = \boxed{\dfrac{\eta_0}{\varepsilon_0}\text{, pointing right, away from the sheets}}
\]
Note $d$ never enters: an infinite sheet's field doesn't depend on distance from it.}

\review{The two outside points have equal field strength, and the field between the
sheets is exactly double that. This is the same pattern as a parallel-plate capacitor
(a much stronger field trapped between the plates than leaks out around them) even
though here the two charge densities aren't equal and opposite $(\eta_0$ vs.\
$3\eta_0)$: the ratio between the ``inside'' and ``outside'' field strength is set purely by
how the two contributions add or cancel, not by the sheets matching each other.}

%============================================================
\problem{23.25}{4}{Figure EX23.25 shows a 1.5 g ball hanging from a string inside a
parallel-plate capacitor made from $12\text{ cm}\times12\text{ cm}$ electrodes, charged
to $\pm75$ nC. The string hangs at $15^\circ$ from vertical, deflected toward the
negative plate. What is the charge on the ball, in nC?}

\model{Treat the capacitor's field as that of an ideal parallel-plate capacitor,
uniform between the plates: $E = Q_{\text{plate}}/(\varepsilon_0 A)$. The ball hangs in
equilibrium under three forces: gravity, string tension, and the electric force
$qE$.}

\visualize{The string is deflected toward the negative plate. The field inside the
capacitor points from the positive plate toward the negative plate. A ball that gets
pushed \emph{along} the field, toward the negative plate, must be positively charged.}

\solve{The plate area is $A = (0.12\text{ m})^2 = 0.0144$ m$^2$, so the field strength is
\[
E = \frac{Q_{\text{plate}}}{\varepsilon_0 A} = \frac{75\times10^{-9}\text{ C}}
{(8.85\times10^{-12})(0.0144\text{ m}^2)} = 5.89\times10^{5}\text{ N/C}.
\]
For the hanging ball, the string's angle from vertical is set by the ratio of the
horizontal electric force to the vertical weight, $\tan\theta = F_{\text{elec}}/(mg)$, so
\[
F_{\text{elec}} = mg\tan\theta = (1.5\times10^{-3}\text{ kg})(9.8\text{ m/s}^2)\tan15^\circ
= 3.94\times10^{-3}\text{ N}.
\]
Then
\[
q = \frac{F_{\text{elec}}}{E} = \frac{3.94\times10^{-3}\text{ N}}{5.89\times10^{5}\text{
N/C}} = \boxed{+6.7\text{ nC}}.
\]}

\review{A $15^\circ$ deflection is fairly gentle, so it makes sense that the electric
force $(3.9\text{ mN})$ comes out smaller than the ball's weight $(mg\approx14.7\text{
mN})$ rather than dwarfing it: $\tan15^\circ\approx0.27$ is a modest ratio, not a
dramatic one. If the deflection had instead been described toward the positive plate,
the same method would give the same magnitude of charge, just negative.}

\end{document}
